% The course readings are written to readings.bib when you compile. Add or edit entries here.
\begin{filecontents*}[overwrite]{readings.bib}
@book{hart2018,
  author    = {Hart, Miriam},
  title     = {Cities and Their Rivers},
  publisher = {Example University Press},
  year      = {2018}
}
@article{mendez2020,
  author  = {Mendez, Carlos},
  title   = {Floodplain planning after 1950},
  journal = {Example Journal of Urban History},
  year    = {2020},
  volume  = {46},
  number  = {2},
  pages   = {233--251}
}
@incollection{ito2017,
  author    = {Ito, Yuki},
  title     = {Embankments as public space},
  booktitle = {Water in the Modern City},
  editor    = {Lang, Oskar},
  publisher = {Sample Books},
  year      = {2017},
  pages     = {77--98}
}
@article{byrne2022,
  author  = {Byrne, Aoife and Chen, Wei},
  title   = {Restoring buried streams: lessons from five projects},
  journal = {Landscape Example},
  year    = {2022},
  volume  = {15},
  pages   = {10--29}
}
@book{fofana2021,
  author    = {Fofana, Aminata},
  title     = {Drought, Memory and Urban Design},
  publisher = {Example Press},
  year      = {2021}
}
\end{filecontents*}
\documentclass[11pt,a4paper]{article}
\usepackage[a4paper,margin=2.2cm]{geometry}
\usepackage[authoryear,round]{natbib}
\usepackage{xcolor}
\usepackage{tabularx}
\usepackage{booktabs}
\usepackage{enumitem}
\usepackage{titlesec}
\usepackage{microtype}
\usepackage[hidelinks]{hyperref}

\definecolor{week}{RGB}{0,95,105}
\setlength{\parindent}{0pt}
\setlength{\parskip}{0.5ex}

\titleformat{\section}{\normalfont\Large\bfseries\color{week}}{}{0pt}{}

% Tags that mark how important a reading is.
\newcommand{\core}{\textsf{\small\colorbox{week}{\textcolor{white}{CORE}}}}
\newcommand{\optional}{\textsf{\small\colorbox{week!15}{\textcolor{week}{OPTIONAL}}}}

% \weekrow{number}{dates}{topic}{readings} adds one row to the schedule table.
\newcommand{\weekrow}[4]{\textbf{#1} & #2 & #3 & #4 \\ \midrule}

\begin{document}

{\small\color{week}\textsf{GEOG 2140 \textbullet\ Autumn 2026}\par}
{\LARGE\bfseries Rivers and the Modern City\par}
{\large Reading list and schedule\par}
\medskip
Convenor: Dr Convenor Name \quad Office hours: Tuesdays 14:00 to 16:00 \quad Room: B2.07

\section{How to use this list}
Read the \core{} items before the seminar in that week. \optional{} items deepen a topic for essays.
Full references are at the end of this document, and each citation links to its entry.

\section{Weekly schedule}
\begin{tabularx}{\linewidth}{@{}l l X X@{}}
\toprule
Week & Dates & Topic & Readings \\
\midrule
\weekrow{1}{22 to 26 Sep}{Why rivers shaped cities}{\core{} \citet[ch.~1]{hart2018}}
\weekrow{2}{29 Sep to 3 Oct}{Planning for floods}{\core{} \citet{mendez2020} \newline \optional{} \citet[ch.~4]{hart2018}}
\weekrow{3}{6 to 10 Oct}{Riverbanks as public space}{\core{} \citet{ito2017}}
\weekrow{4}{13 to 17 Oct}{Burying and restoring streams}{\core{} \citet{byrne2022}}
\weekrow{5}{20 to 24 Oct}{Reading week}{No new readings. Start your essay plan.}
\weekrow{6}{27 to 31 Oct}{Scarcity and drought}{\core{} \citet[chs.~2 and~3]{fofana2021}}
\end{tabularx}

\section{Assessment}
\begin{itemize}[nosep]
  \item Seminar reflections, 20\%: due each Friday of weeks 2 to 6.
  \item Essay of 2\,500 words, 80\%: due 15 December 2026.
\end{itemize}

\section{Further reading}
Add sources for students who want to go further. \nocite{*}The command on this line lists every entry in the file, even ones not cited above.

\renewcommand{\refname}{Full references}
\bibliographystyle{plainnat}
\bibliography{readings}

\end{document}
